% ======================================================================
%  Course 2610EST3000 -- Project 2 (Autoresearch-Base) report template
%  Typeset in the official NeurIPS 2025 format, preprint (non-anonymous)
%  variant, with course information in the running header/footer.
% ======================================================================
%  Build:  make            (pdflatex x2 + bibtex)
%  Clean:  make clean
%
%  This is a COURSE submission, not a conference submission:
%  * Use the preprint option so author names (with student IDs) are shown
%    and no "Submitted to NeurIPS" notice or line numbers appear.
%  * The first-page bottom notice is repurposed as the course notice.
%  * Pages 2+ carry a running header (course ID / project) and a footer
%    with the Canvas link and page numbers.
%  * The NeurIPS Paper Checklist stays: treat it as the self-review
%    rubric for your report -- answer honestly.
%  * Never modify neurips_2025.sty and never \renewcommand margins, font
%    sizes, or line spacing.
%
%  Line convention for this template: one paragraph per line, no hard
%  wrapping. Let your editor soft-wrap.
% ======================================================================
\documentclass{article}

% --- Course submission version (non-anonymous, authors + student IDs shown) ---
\usepackage[preprint]{neurips_2025}

\usepackage[utf8]{inputenc}
\usepackage[T1]{fontenc}
\usepackage{hyperref}       % cross-references and links
\usepackage{url}
\usepackage{booktabs}       % the only correct table style at AI venues (\toprule/\midrule/\bottomrule)
\usepackage{amsfonts}
\usepackage{amsmath}
\usepackage{amssymb}
\usepackage{nicefrac}
\usepackage{microtype}
\usepackage{xcolor}
\usepackage{graphicx}       % vector PDF figures only: \includegraphics{fig.pdf}
\usepackage{subcaption}     % subfigures (a)(b)(c)
\usepackage{algorithm}
\usepackage{algpseudocode}  % pseudocode floats
\usepackage{fancyhdr}       % running header/footer with course info
\usepackage{lastpage}       % "Page X of Y"

% ----------------------------------------------------------------------
%  Course identity -- edit COURSENAME only if the instructor says so
% ----------------------------------------------------------------------
\newcommand{\courseid}{2610EST3000}
\newcommand{\courseproject}{Project 2: Recursive Self-Improvement Strategies}
\newcommand{\coursecanvas}{https://canvas.westlake.edu.cn/courses/2063}
\definecolor{coursegray}{gray}{0.45}

% Repurpose the first-page bottom notice (in preprint mode it would print
% "Preprint.") as the course notice -- the natural home for it.
\makeatletter
\renewcommand{\@noticestring}{%
  \courseid\ \ $\cdot$ \ \ \courseproject\\[1pt]
  Submitted via Canvas: \url{\coursecanvas}%
}
\makeatother

% Running header/footer for pages 2+ (page 1 is kept clean by the style).
\pagestyle{fancy}
\fancyhf{}
\renewcommand{\headrulewidth}{0.3pt}
\renewcommand{\footrulewidth}{0pt}
\fancyhead[L]{\footnotesize\color{coursegray}\courseid\ \ $\cdot$ \ \ \courseproject}
\fancyhead[R]{\footnotesize\color{coursegray}Page \thepage\ of \pageref{LastPage}}
\fancyfoot[C]{\footnotesize\color{coursegray}\url{\coursecanvas}}

% ----------------------------------------------------------------------
%  Writing aids (must be zeroed out before submission)
% ----------------------------------------------------------------------
\newcommand{\todo}[1]{\textcolor{red}{\textbf{[TODO: #1]}}}   % grep for TODO before submitting; none may remain
\newcommand{\cm}[1]{\textcolor{blue}{[comment: #1]}}           % collaborator margin notes

% ======================================================================
\title{\todo{Title: state the strategy you built and its measurable payoff in one line; avoid "A Study of" titles}}
%  Good shape:  "Evidence-Guided Proposal and Guarded Self-Edits Improve
%               the Yield of a Recursive Self-Improvement Engine"
%  Bad shapes:  "Research on Self-Improving Systems" (vague),
%               "A Novel Framework for ..." (cliche)

% Course submissions are NON-anonymous: every author's full name, student
% ID, and Westlake email must appear. Only one member is allowed. 
\author{%
  Firstname Lastname\\
  Student ID: \todo{2024XXXXXX}\\
  Westlake University\\
  \texttt{firstname.lastname@westlake.edu.cn}
}

\begin{document}

\hypersetup{pageanchor=false}  % paper restarts at page 1 below; avoid duplicate anchors
\pagenumbering{roman}          % cover page is "i" (never displayed)
\thispagestyle{empty}

% This cover page alone uses 1cm margins so the rubric can breathe;
% the paper's own geometry is restored verbatim below.
\newgeometry{left=1cm,right=1cm,top=1cm,bottom=1cm}

\begin{center}
{\LARGE\bf Project 2 Assessment Guide\par}
\vspace{8pt}
{\large \courseid\ \ $\cdot$ \ \ Project 2: Recursive Self-Improvement Strategies (100 pts) \par}
\end{center}

\vspace{2pt}
\noindent\rule{\textwidth}{1.2pt}\\[8pt]
\noindent \textcolor{red}{This cover page states how the report that follows is graded. It is an instructor annotation: it is not part of the report, does not count toward the 9-page body budget, and the report begins on the next page with its usual layout and page numbering.}

% Project Introduction.
\vspace{8pt}
\noindent{\normalsize\bf About the project}\\[4pt]
\noindent{\normalsize\textbf{Background.} A recursive self-improvement (RSI) engine runs research cycles autonomously: each cycle proposes a single-variable experiment, runs it under a fixed budget, records the outcome, and may rewrite the engine's own strategy through a guard that admits only edits with immediate compute savings.}\\[2pt]
\noindent{\normalsize\textbf{Problem.} Naive strategies waste most cycles on re-exploration and no-change experiments, so accumulated knowledge does not raise the research yield and the engine plateaus early.}\\[2pt]
\noindent{\normalsize\textbf{Requirements.} Without editing the frozen engine, design, motivate, implement and evaluate strategies for proposal generation, result diagnosis and guarded self-editing; report evidence against matched baselines over fixed seeds, with the fate of every proposed self-edit accounted for.}\\[2pt]
\noindent{\normalsize\textbf{Learning outcomes.} Formulate a research bottleneck as a learning problem; connect the design to course concepts; run controlled experiments with honest statistics; and communicate results with evidence-calibrated claims.}

\vspace{10pt}
\begin{center}
\small
\begin{tabular}{@{}p{3.1cm}p{1.4cm}p{14.1cm}@{}}
\toprule
\textbf{Component} & \textbf{Weight} & \textbf{What earns the marks} \\
\midrule
Literature survey & 10\% & \textbf{Coverage:} the foundations of recursive self-improvement and automated research, plus the most relevant recent work. \newline \textbf{Precision:} each citation supports the exact claim it is attached to; no padding. \newline \textbf{Synthesis:} approaches are grouped, compared and contrasted, not listed. \newline \textbf{Gap:} the limitation this submission targets is named, with evidence from the literature. \newline \textbf{Positioning:} what is reused from the given engine is separated from what is new here. \\
\addlinespace[5pt]
\midrule
Reasoning \& insight & 30\% & \textbf{Diagnosis:} the bottleneck (e.g., wasted cycles, no-op edits) is identified from measured behaviour, not intuition. \newline \textbf{Motivation:} every strategy states the hypothesis it tests and why its mechanism should close the gap. \newline \textbf{Course linkage:} explicit use of ML concepts --- exploration--exploitation, bias--variance, credit assignment, MDP/bandit formulation, generalization, statistical testing. \newline \textbf{Formalization:} the decision problem (what to propose, keep, or edit) is posed with assumptions stated. \newline \textbf{Theory:} derivations, propositions, or a compute/budget argument from which the design follows. \newline \textbf{Alternatives:} rejected designs and reasons; predictions that the experiments can falsify. \\
\addlinespace[5pt]
\midrule
Experiments & 30\% & \textbf{Implementation:} every claimed strategy exists in runnable code, integrated within the frozen engine and trust boundary. \newline \textbf{Protocol:} naive and matched baselines, equal cycle budgets, multiple seeds fixed in advance, metric defined before the runs. \newline \textbf{Main results:} seed mean $\pm$ standard deviation against baselines; effect sizes quantified. \newline \textbf{Offline evaluation (anti-overfitting):} to prevent overfitting to the public benchmarks, the course staff will evaluate all submissions on five unseen offline (held-out) tasks; hardcoded shortcuts will be penalized. \newline \textbf{Reproducibility:} exact commands, configs, seed lists and the self-edit ledger; a clean checkout reproduces the headline number. \newline \textbf{Honesty:} negative and failed runs reported; compute cost accounted; no cherry-picking. \\
\addlinespace[5pt]
\midrule
Result analysis & 30\% & \textbf{Attribution:} ablations isolate which component carries the gain. \newline \textbf{Mechanism:} gains and failures explained causally --- overfitting to noise, distribution shift, non-stationarity, exploration collapse. \newline \textbf{Dynamics:} keep-rate versus budget, per-task and per-cycle behaviour, where gains and stagnation arise. \newline \textbf{Failure anatomy:} rejected proposals, discarded self-edits and guard arbitrations are taxonomized. \newline \textbf{Calibration:} limitations and threats to validity stated; conclusions match the strength of the evidence. \\
\bottomrule
\end{tabular}
\end{center}

\vspace{6pt}
\noindent{\small\textbf{Marking.} Component marks are weighted by the percentages above and summed to 100; within a component, credit follows the strength of verifiable evidence, from fully met (85--100\%) through met but thin (70--84\%) and partial or asserted (60--69\%) to unsupported ($<$60\%).}

\vspace{4pt}
\noindent{\small\textbf{Deliverables (Canvas).} The report PDF (this template; at most 9 pages of body, excluding references, the checklist and the appendix) and a runnable code archive implementing every reported strategy, whose clean checkout reproduces the headline number, together with seed lists, run logs and the self-edit ledger.}

\vspace{4pt}
\noindent{\small\textbf{Policies.} Individual work, one author per submission; the frozen engine and trust boundary may not be modified; every claim must trace to actual runs and real, verifiable citations --- fabricated results or references score zero; LLM use in methods and in writing disclosed per the checklist; deadline and late policy per the Canvas announcement.}

\noindent\rule{\textwidth}{1.2pt}\\[8pt]

\clearpage
% Restore the paper's geometry exactly as neurips_2025.sty sets it.
\newgeometry{textheight=9in,textwidth=5.5in,top=1in,headheight=12pt,headsep=25pt,footskip=30pt}
\pagenumbering{arabic}         % paper restarts at 1: original layout and page count
\hypersetup{pageanchor=true}
% ======================================================================

\maketitle

% ======================================================================
%  Body (total budget 9 pages; per-section budgets in each file header)
% ======================================================================
% ======================================================================
%  ABSTRACT -- budget: 150-250 words, one paragraph, no citations
% ======================================================================
%  The de-facto four-move skeleton at AI venues:
%   1. Problem: the task and why it is hard, in one sentence (no background warm-up).
%   2. Insight: the key idea ("We observe that ..." / "Our key insight is ...").
%   3. Method: what you actually built (strategy names + mechanism, not a pipeline recap).
%   4. Result: the single strongest quantified claim, with concrete numbers and the baseline.
%
%  Taboos: opening with "In this paper, we ..."; no numbers; claiming more than the experiments support.
% ======================================================================
\begin{abstract}
Recursive self-improvement (RSI) engines can run autonomously, but with naive strategies their research yield plateaus quickly. \todo{1-2 sentences: name the specific bottleneck you attack, e.g. "proposals ignore accumulated evidence, so most cycles are wasted on re-exploration".} Our key insight is that \todo{one-sentence core insight}. We implement \todo{the strategy bundle, corresponding to whichever of C-1/C-2/E-2/E-3 you completed} on top of a frozen three-loop RSI engine, and evaluate on \todo{task name(s)} under matched cycle budgets. \todo{Method name} improves the final best result from \todo{X} to \todo{Y} ($\pm$ \todo{seed std} over \todo{N} seeds), raises proposal keep-rate from \todo{A\%} to \todo{B\%}, and produces \todo{Z} promoted self-edits that survive guard arbitration. Ablations show that \todo{which component carries the gain}.
\end{abstract}
        % ~0.25 column
% ======================================================================
%  Section 1 INTRODUCTION -- budget: ~1 page
% ======================================================================
%  Standard six-paragraph logic chain (one function per paragraph, fixed order):
%   P1 Context + running example: open from a concrete scenario; one sentence that places the problem domain.
%   P2 Limitations of the status quo: exactly how the naive strategies fail (cite the reference experiment's plateau as evidence).
%   P3 Problem essence + this paper's goal: formalize in one sentence.
%   P4 Key challenges: why this is not "just tuning" (trust boundary / measurement noise / causal attribution).
%   P5 Solution overview: how each strategy module answers a challenge (one-to-one mapping).
%   P6 Contributions: 3-4 itemized, verifiable claims; no adjectives.
%
%  The reviewer's first impression forms here: paragraph one must be concrete, and every contribution must be cashed out later in the paper.
%  Taboos: implementation details in the intro; writing "we propose a novel framework" as a contribution.
% ======================================================================
\section{Introduction}
\label{sec:intro}

% --- P1: context and running example ---
\todo{Open with the RSI engine running a research task: the single-variable experiment loop, the accumulating knowledge base, self-edits that raise research yield. Two or three sentences of minimal necessary background; assume an ML-literate reader who has never heard of this system.}

% --- P2: status quo and its failure modes ---
\todo{Be concrete about the naive failures: the fixed grid walk ignores evidence (cite your own measurements on the public tasks: what fraction of cycles were no-change waste; after which cycle the best result plateaus; where it falls short of the Claude Code baseline in \texttt{baselines/public\_tasks\_baseline.md}).}

% --- P3: problem and goal ---
\todo{One sentence: "with the engine loop and trust boundary, and the five public task scorers, frozen, make the strategy layer open a reproducible yield gap over the naive baseline and close the gap to the Claude Code baseline."}

% --- P4: challenges ---
\todo{At least two: (i) causal attribution of keep/discard decisions under observation noise; (ii) self-edits must pass the guard without breaking determinism; add the ones you actually hit.}

% --- P5: solution overview ---
\todo{One sentence per challenge mapping to a strategy module, with a forward reference to \S\ref{sec:method}.}

% --- P6: contributions ---
Our contributions are:
\begin{itemize}
  \item \todo{Contribution 1: a strategy plus its one-line mechanism (maps to C-1/E-3/...).}
  \item \todo{Contribution 2: same pattern.}
  \item \todo{Contribution 3: an empirical finding with numbers, e.g. "evidence-guided proposal doubles keep-rate at equal budget (Table~\ref{tab:main})".}
  \item \todo{(Optional) a negative result or analysis contribution; negative results are respected at AI venues.}
\end{itemize}
    % ~1 page
% ======================================================================
%  Section 2 RELATED WORK -- budget: 0.5-0.75 page
% ======================================================================
%  Organize by theme, one paragraph each, and end every paragraph with a single sentence stating how we differ.
%  Three natural threads for this project:
%   (a) Self-referential / self-improving systems: G\"odel machine \citep{schmidhuber2009godel}, STOP \citep{zelikman2023stop}, Darwin-G\"odel Machine \citep{zhang2025dgm}.
%   (b) LLM research-automation agents: AI Scientist \citep{lu2024aiscientist}, MLAgentBench \citep{huang2024mlagentbench}, RE-Bench \citep{wijk2024rebench}.
%   (c) Automated algorithm discovery / optimization: AutoML-Zero \citep{real2020automlzero}, OPRO \citep{yang2024opro}, AlphaEvolve \citep{novikov2025alphaevolve}.
%
%  Conventions: high citation density (3-8 per paragraph); \citet when authors are the subject, \citep for parenthetical use; cite only what you have actually read, and every citation must be used by a sentence (no padding).
%  The closing positioning paragraph is mandatory ("Ours differs in that ..."): related work without positioning is a reference list, and reviewers deduct for it.
% ======================================================================
\section{Related Work}
\label{sec:related}

\paragraph{Self-referential and self-improving systems.}
\todo{From the theoretical origin \citep{schmidhuber2009godel} to LLM-era self-rewriting \citep{zelikman2023stop,zhang2025dgm}. Point out: those works rewrite \emph{task-solving code}; we rewrite the \emph{research strategy itself}, under a guard.}

\paragraph{LLM agents for research automation.}
\todo{\citep{lu2024aiscientist,huang2024mlagentbench,wijk2024rebench}: the state of end-to-end research agents and their evaluation. Point out: they use the LLM as the engine; we freeze the engine and make the research strategy itself the replaceable, measurable object.}

\paragraph{Automated algorithm discovery.}
\todo{\citep{real2020automlzero,yang2024opro,novikov2025alphaevolve}: searched/evolved algorithms often beat hand design but demand large compute and human-in-the-loop evaluation. Point out: we study the same question under CPU-second, deterministic, reproducible constraints.}

\paragraph{Positioning.}
\todo{Two or three closing sentences: relative to everything above, what is the unique coordinate of this work (frozen engine + strategy-layer science + strict keep/discard causal measurement)?}
    % ~0.5-0.75 page
% ======================================================================
%  Section 3 THE BASE SYSTEM -- budget: ~0.75 page
% ======================================================================
%  Purpose: give the reader just enough to understand which seams your strategies plug into in Section 4 -- nothing more.
%  This is the given machinery, not your contribution; writing too much here dilutes the contribution signal.
%  Three things suffice:
%   1. The three loops (fast/middle/slow) and what each does; one small figure or one tight paragraph.
%   2. The frozen contracts and the trust boundary (EDITABLE\_SURFACES, guard layers, keep|discard|fail).
%   3. The strategy-seam table (Protocols: Proposer / SelfImprover / Gates / Diagnostician / ...), one line per seam with a pointer to the Section 4 subsection that fills it.
%
%  A system figure is strongly recommended (Figure 1): the three loops plus the self-edit circuit.
%  Figure conventions: vector PDF; in-figure text no smaller than footnote size (~7pt at final size); the caption must stand alone (understandable without the body text).
% ======================================================================
\section{The Base System}
\label{sec:system}

\todo{Two or three sentences: given a task bundle (one editable file + one gate metric + a fixed budget), the engine runs three nested loops; reference Figure~\ref{fig:system}.}

\paragraph{Loops.}
\todo{Fast: propose $\to$ implement $\to$ run $\to$ finalize $\to$ analyze, single-variable experiments. Middle: the benchmark is frozen within an epoch and upgraded across epochs; self-edits are arbitrated here. Slow: hypotheses succeed each other on a typed-edge question graph; saturated questions spawn children.}

\paragraph{Frozen contracts and the trust boundary.}
\todo{The concept contracts, the verdict trichotomy, the guard layers; emphasize that self-edits may only land in the sentinel region of \texttt{state/strategy.py}, under the applying $\to$ observing $\to$ kept|reverted lifecycle with immediate re-entry.}

\paragraph{Strategy seams.}
\todo{A small table: seam name / Protocol / naive default behavior / where your implementation lives (section pointer). This table is the bridge from Section 3 to Section 4 and reviewers respond well to it.}

% Placeholder figure: replace the \fbox with
%   \includegraphics[width=\linewidth]{figures/system.pdf}
% once the diagram exists. Note the label must be INSIDE \caption's float.
\begin{figure}[t]
  \centering
  \fbox{\parbox[c][5cm][c]{0.9\linewidth}{\centering \todo{System diagram: three loops + self-edit circuit; vector PDF; frozen components grey, yours in color.}}}
  \caption{The three loops and the self-edit circuit. \todo{Make the caption complete: one clause per loop, and mark which components are frozen (grey) versus implemented in this work (color).}}
  \label{fig:system}
\end{figure}
          % ~0.75 page (the given machinery -- keep brief)
% ======================================================================
%  Section 4 METHOD -- budget: 1.5-2 pages; the core of the paper
% ======================================================================
%  One subsection per strategy you actually completed (e.g.):
%   4.1 Evidence-Guided Proposer (C-1)
%   4.2 Waste-Triggered Self-Improver (C-2)
%   4.3 Noise-Floor Diagnostics & Gates (E-3/E-2)
%  Inside each subsection, use a uniform three-part pattern:
%   (i)   Motivation: the concrete failure of the naive version plus a one-sentence insight.
%   (ii)  Mechanism: a precise algorithmic description -- pseudocode (algorithm float) or equations; write to "a reader could re-implement this" granularity; never "we use a model to ..." hand-waving.
%   (iii) Constraints honored: why the design satisfies the frozen contracts (determinism, guard, Protocol signatures).
%
%  Conventions:
%  * Keep notation consistent across the whole section and define every symbol at first use.
%  * Pseudocode earns more trust than prose -- give the core strategy an Algorithm float.
%  * For every design choice, answer "why not the simpler alternative" -- this pre-empts reviewer questions.
% ======================================================================
\section{Method: Strategies Behind the Seams}
\label{sec:method}

\todo{Open with two or three overview sentences: which strategies you implemented and how they cooperate (e.g. the proposer generates candidates, diagnostics decides when to upgrade the epoch, the improver turns waste signals into strategy edits). Then one subsection per strategy.}

\subsection{\todo{Strategy name, e.g. Evidence-Guided Proposal Ranking} (C-1)}
\label{sec:method:proposer}

\paragraph{Motivation.}
\todo{The concrete naive failure: fixed grid, ignores KB evidence $\to$ cite your measurement (e.g. what fraction of the naive run's 292 cycles were repeated or wasted proposals).}

\paragraph{Mechanism.}
\todo{Precisely: how candidates are ranked from \texttt{evidence\_for} / \texttt{best\_result} / \texttt{uncertainty}; how the value grid adapts. Provide an equation or pseudocode:}

% \begin{algorithm}[t]
% \caption{\todo{Strategy name}}
% \label{alg:proposer}
% \begin{algorithmic}[1]
%   \Require knowledge base $\mathcal{K}$, frontier question $q$, budget $B$
%   \State \todo{...}
% \end{algorithmic}
% \end{algorithm}

\paragraph{Constraints honored.}
\todo{Determinism (the seeded RNG as the only randomness source), unchanged Protocol signatures, no touching frozen components -- one sentence each.}

\subsection{\todo{Strategy name} (C-2 / E-3 / ...)}
\label{sec:method:second}
\todo{Follow the same three-part pattern.}
          % ~1.5-2 pages (your strategies -- the core)
% ======================================================================
%  Section 5 EXPERIMENTS -- budget: ~1.5 pages (setup + main results)
% ======================================================================
%  Structure: 5.1 Setup, 5.2 Main results. Anything analytical goes to Section 6.
%
%  Setup checklist (this is the reviewer's reproducibility scorecard):
%   * Task: which bundle, the metric definition, the budget (cycles / wall-clock).
%   * Baselines: the given naive one plus a noop control; every strategy needs at least one "switch it off" ablation slot.
%   * Seeds: N >= 3, report mean +/- std; single-seed results are not accepted at AI venues.
%   * Fairness: compare at equal cycle budget; state any asymmetry explicitly.
%   * Hyperparameters: list them all (body or appendix), including how they were chosen (never tune on the test metric).
%
%  Result presentation conventions:
%   * Main table (Table 1): rows = methods, columns = metrics; booktabs three-line style, no vertical rules; bold the best entry only when its gap over the runner-up is significant.
%   * Main figure: trajectory curves must carry a seed-variance band (mean +/- std/sem shading).
%   * Every number must be traceable: which experiment directory and which command produced it (full details in the appendix).
%
%  Honesty red line: a keep-rate increase with a flat best result is NOT a yield improvement; never let a secondary metric's improvement paper over a plateau on the primary one. ANALYSIS.md section 3 of the reference experiment is a model of this kind of honest analysis.
% ======================================================================
\section{Experiments}
\label{sec:exp}

\subsection{Setup}
\label{sec:exp:setup}

\paragraph{Tasks and metrics.}
\todo{The five public tasks (\texttt{battle-of-sexes}, \texttt{prisoners-dilemma}, \texttt{colonel-blotto}, \texttt{house-price}, \texttt{breakout-lite}): one-line task definition + the gate metric (higher/lower is better) + where the headroom comes from. Development tasks \texttt{quad-optim}/\texttt{iris-ood} may be mentioned but do not replace the five.}

\paragraph{Baselines.}
\todo{The given naive strategy, the noop-improver control, the Claude Code strong baseline (\texttt{baselines/public\_tasks\_baseline.md}), and one ablation slot per strategy of yours. Explain why these are fair opponents.}

\paragraph{Protocol.}
\todo{Cycles, seeds (N=\todo{}), the determinism guarantee, hardware (CPU model suffices), wall-clock per run. Command-level reproducibility goes to Appendix~\ref{app:repro}.}

\subsection{Main Results}
\label{sec:exp:main}

\todo{Lead with the conclusion sentence for Table~\ref{tab:main} ("X outperforms naive by ..."), then interpret column by column; never just drop a table. Every claim must cite a specific table/figure number.}

\begin{table}[t]
  \caption{Main results: final best result and keep-rate, mean $\pm$ std over $N$ seeds. \todo{State budget and seed count in the caption; the table must be self-contained. Bold the best entry only when its gap over the runner-up is significant.}}
  \label{tab:main}
  \centering
  \begin{tabular}{lcc}
    \toprule
    Method & Final best $\uparrow$ & Keep-rate $\uparrow$ \\
    \midrule
    Naive (given)        & \todo{} & \todo{} \\
    Noop-improver ctrl.  & \todo{} & \todo{} \\
    Ours (self-edit ON)  & \todo{} & \todo{} \\
    Ours (self-edit OFF) & \todo{} & \todo{} \\
    Claude Code baseline & \todo{} & --- \\
    \bottomrule
  \end{tabular}
\end{table}
     % ~1.5 pages (setup + main results)
% ======================================================================
%  Section 6 ANALYSIS -- budget: 1.5-2 pages; as important as Section 5
% ======================================================================
%  The main results show THAT it works; this section answers WHY it works and WHEN it fails.
%  The assignment's core narrative also lives here: the end-to-end causal story of at least one promoted self-edit (what changed in state/strategy.py, and which measured quantity improved after re-entry).
%
%  Suggested subsections (keep whichever you have material for):
%   6.1 Ablation: how much each component loses when removed -- ablation is the only evidence attributing the gain.
%   6.2 Anatomy of a promoted self-edit: waste signal $\to$ guard passage $\to$ observing-window measurement $\to$ kept. Pair with a before/after trajectory figure.
%   6.3 Where it fails: negative results and failure modes. Writing these honestly earns credit.
%   6.4 Cost: wall-clock / per-cycle overhead, proving the strategy is not trading compute for yield.
%
%  Conventions: open each analysis subsection with a one-sentence research question and close it with a one-sentence answer. Diagnostic figures (keep-rate curves, edit-response plots, question-graph snapshots) must be vector PDFs.
% ======================================================================
\section{Results}
\label{sec:analysis}

\subsection{Performance}
\label{sec:analysis:performance}
\todo{Research question: what is the overall performance across tasks? Show the main metric (e.g., task scores), with error bars. One-sentence answer: the strategy improves yield by X\% over baseline.}

\subsection{Anatomy of a Promoted Self-Edit}
\label{sec:analysis:selfedit}
\todo{The assignment's core requirement, told as a timeline: which waste pattern appeared at cycle N $\to$ which edit the improver proposed (quote the key before/after fragment of \texttt{state/strategy.py}) $\to$ each guard layer's verdict $\to$ the M-cycle observing-window comparison $\to$ the kept/reverted ruling. Pair it with a trajectory figure that marks the promotion moment with a vertical line.}

\subsection{Failure Modes}
\label{sec:analysis:failure}
\todo{Document at least one real failure: improver proposals that get reverted, noise-induced false keeps, question-graph evolution walking into a dead end. If you omit negative results, reviewers will ask.}

\subsection{Cost-Effectiveness}
\label{sec:analysis:cost}
\todo{Extra wall-clock per cycle, token spend (if the LLM seam is used), and whether determinism is preserved.}

\subsection{Module Ablations}
\label{sec:analysis:ablation}
\todo{Research question: which component produces the main gain? Table: full system versus one component removed at a time. When interpreting, distinguish "significant drop" from "within noise".}
        % ~1.5-2 pages (ablations / anatomy / failures)
% ======================================================================
%  Section 7 DISCUSSION -- budget: ~0.5 page
% ======================================================================
%  Two blocks, neither may be vacuous:
%  * Limitations: the real boundaries (single task domain, CPU budget, LLM seam unused, cross-epoch comparability), each with one sentence on how it could be lifted. Writing "we leave X to future work" without admitting a present weakness loses points.
%  * Broader impact: the dual-use nature of self-improving systems -- your guard / meltdown brake safety designs are exactly the positive argument; state them honestly.
% ======================================================================
\section{Discussion}
\label{sec:discussion}

\paragraph{Limitations.}
\todo{Two or three genuine limitations, each with one sentence on the boundary and one on a path forward.}

\paragraph{Broader impact.}
\todo{The risks of self-improving engines (uncontrolled self-modification, metric gaming) and this work's corresponding designs (trust boundary, default-reject guard, meltdown brake, deterministic auditing).}
      % ~0.5 page (limitations + broader impact)
% ======================================================================
%  Section 8 CONCLUSION -- budget: ~0.25 page, 3-4 sentences, no new information
% ======================================================================
%  Conventions: restate the problem + one sentence of method + the single strongest quantified result + one open direction. Introduce no claim absent from the body; do not replay the abstract; do not open with "In conclusion".
% ======================================================================
\section{Conclusion}
\label{sec:conclusion}

\todo{Three or four closing sentences.}
      % ~0.25 page

% ======================================================================
%  References -- page limit does not apply from here on.
%  Convention: natbib author-year; use \citep{} for parenthetical
%  citations and \citet{} when the authors are the sentence subject;
%  never hand-write [1] [2].
% ======================================================================
\bibliographystyle{plainnat}
\bibliography{references}

% ======================================================================
%  NeurIPS Paper Checklist -- mandatory; omission is grounds for rejection
% ======================================================================
% ======================================================================
%  NeurIPS Paper Checklist -- mandatory; follows the references; does not count toward the page limit
% ======================================================================
%  Rules: every item takes exactly one answer, [Yes] / [No] / [NA], plus a short justification with a section pointer.
%  A "No" answer does not automatically reject the paper, but a missing or perfunctory checklist does.
%  All answers must be truthful -- a fabricated checklist is an ethics violation.
%  This skeleton follows the NeurIPS 2025 checklist structure; re-check against the official current-year version before submitting.
% ======================================================================
\section*{NeurIPS Paper Checklist}

\begin{enumerate}

\item {\bf Claims}
\item[] Question: Do the main claims made in the abstract and introduction accurately reflect the paper's contributions and scope?
\item[] Answer: \todo{[Yes/No/NA]}
\item[] Justification: \todo{Point to the contribution list in \S\ref{sec:intro} and the corresponding evidence in \S\ref{sec:exp}.}

\item {\bf Limitations}
\item[] Question: Does the paper discuss the limitations of the work performed by the authors?
\item[] Answer: \todo{[Yes]}
\item[] Justification: \todo{Point to \S\ref{sec:discussion}.}

\item {\bf Theory assumptions and proofs}
\item[] Question: For each theoretical result, does the paper provide the full set of assumptions and a complete (and correct) proof?
\item[] Answer: \todo{[NA] -- this is empirical work; if you make any theoretical claim, point to its proof here.}

\item {\bf Experimental result reproducibility}
\item[] Question: Does the paper fully disclose all the information needed to reproduce the main experimental results?
\item[] Answer: \todo{[Yes]}
\item[] Justification: \todo{Point to Appendix~\ref{app:repro}: commands, seeds, configuration. The system is deterministic -- same seed, same result.}

\item {\bf Open access to data and code}
\item[] Question: Does the paper provide open access to the data and code?
\item[] Answer: \todo{[Yes/No]}
\item[] Justification: \todo{State where the code and data live (course repository or public link). This is a course submission -- no anonymization is needed.}

\item {\bf Experimental setting/details}
\item[] Question: Does the paper specify all the training and test details?
\item[] Answer: \todo{[Yes]}
\item[] Justification: \todo{Point to \S\ref{sec:exp:setup} and the appendix hyperparameter table.}

\item {\bf Experiment statistical significance}
\item[] Question: Does the paper report error bars suitably and correctly defined or other appropriate information about statistical significance?
\item[] Answer: \todo{[Yes]}
\item[] Justification: \todo{Mean $\pm$ std over N seeds; name the significance test used.}

\item {\bf Experiments compute resources}
\item[] Question: For each experiment, does the paper provide sufficient information on the computer resources needed?
\item[] Answer: \todo{[Yes]}
\item[] Justification: \todo{CPU model, wall-clock per run, total compute hours.}

\item {\bf Code of ethics}
\item[] Question: Does the research conform to the NeurIPS Code of Ethics?
\item[] Answer: \todo{[Yes]}

\item {\bf Broader impacts}
\item[] Question: Does the paper discuss both potential positive societal impacts and negative societal impacts?
\item[] Answer: \todo{[Yes]}
\item[] Justification: \todo{Point to the broader-impact paragraph in \S\ref{sec:discussion}.}

\item {\bf Safeguards}
\item[] Question: Does the paper describe safeguards for responsible release of high-risk models or data?
\item[] Answer: \todo{[Yes/NA]}
\item[] Justification: \todo{The guard, trust boundary, and meltdown brake are the safeguards; point to \S\ref{sec:system} and \S\ref{sec:discussion}.}

\item {\bf Licenses for existing assets}
\item[] Question: Are the creators or original owners of assets (code, data, models) properly credited, and are licenses and terms of use respected?
\item[] Answer: \todo{[Yes]}

\item {\bf New assets}
\item[] Question: Are new assets introduced in this paper well documented, and is the documentation alongside the assets?
\item[] Answer: \todo{[Yes/NA]}

\item {\bf Crowdsourcing and research with human subjects}
\item[] Answer: \todo{[NA]}

\item {\bf Institutional review board (IRB) approvals}
\item[] Answer: \todo{[NA]}

\item {\bf Declaration of LLM usage}
\item[] Question: Does the paper describe the usage of LLMs if it is an important, original, or non-standard component of the core methods?
\item[] Answer: \todo{[Yes/NA]}
\item[] Justification: \todo{If your strategies use the \texttt{rsi.llm} seam, state the model, the role of each prompt, and how caching preserves determinism; also disclose any LLM assistance in writing (NeurIPS requires disclosure, and LLMs cannot be authors).}

\end{enumerate}


% ======================================================================
%  Appendix -- does not count toward the page limit, but reviewers are
%  not obliged to read it; every claim supporting a core result must
%  also appear (in summary form) in the body.
% ======================================================================
\appendix
% ======================================================================
%  Appendix: Reproducibility -- the trust foundation reviewers check first
% ======================================================================
%  Goal: a reviewer who does not know you can re-produce the main table within ten minutes using only this appendix.
%  The system's determinism (same seed, same result) is a strength -- state it as one.
% ======================================================================
\section{Reproducibility Details}
\label{app:repro}

\paragraph{Environment.}
\todo{OS, Python version, \texttt{environment.yml}; CPU model and memory.}

\paragraph{Exact commands.}
\todo{The full command behind every row of the main table, for example:}
% \begin{verbatim}
% python -m rsi.core --task tasks/iris-ood --out out/main --cycles 292 --seed 7
% python scripts/report.py --run out/main --out report/
% \end{verbatim}

\paragraph{Seeds.}
\todo{The full seed list; state that seeds were fixed in advance, not selected post hoc.}

\paragraph{Configuration.}
\todo{The full config.yaml or its diff from default; every non-default hyperparameter.}

\paragraph{Self-edit ledger.}
\todo{The complete ledger of promoted and reverted edits (terminal states), backing the causal story in \S\ref{sec:analysis:selfedit}.}

% ======================================================================
\section{Extended Results}
\label{app:extended}
\todo{What did not fit in the body: per-seed trajectories, additional ablations, question-graph end states, the full failure-case catalog. Reviewers are not obliged to read the appendix -- every core conclusion must stand in the body alone.}


\end{document}
